\documentclass[12pt,        % standard font size
  english,ngerman,          % english primary, german secondary
  paper=a4,                 % standard paper size
  captions=tablesignature,  % captions below tables
  listof=numbered,          % including lists of... in the ToC
  bibliography=totoc,       % references in ToC
  headings=small,           % size of headings
  headinclude=false,        % don't include page head in page layout
  footinclude=false,        % don't include page foot in page layout
  parskip=half-,            % space between paragraphs, no indentation
%   dotlessnumbers,         % no dots after chapter etc. numbers
% FIXME Adjust following 2 options for two-sided print, nr of pages
  oneside,                  % one-sided print
%   twoside,                % two-sided print
% FIXME Adjust the Binding CORrection to the binding and number of pages
% 15mm should be fine for theses with about 60 to 120 double sided pages
% and with the standard binding for theses (hardcover)
% Add 2-3mm for longer theses, substract 2-3 mm for shorter theses
% For more information see the documentation for KOMA-classes
%  BCOR=15mm,                 % BCOR for thesis (hardcover)
%  BCOR=5mm,                 % BCOR for project reports (plastic folder)
%                             % No BCOR for seminar reports (plastic folder)
%   draft,                  % speed things up before final version
%   DIV=calc                % calculate the page layout
  DIV=12                    %  
%   DIV=classic             %
  ]{scrbook}                % KOMA script book class

% most setup is done in a file of its own
% this makes ``References'' a section instead of a chapter
% this has to be done before other packages mess around
% with the definition
\makeatletter
\renewcommand*\bib@heading{%
	\section*{\bibname}%
	\@mkboth{\bibname}{\bibname}%
}

% Change colour of superscriptes numbers from footnotes
\renewcommand\@makefnmark{\hbox{
	\@textsuperscript{\normalfont\color{coliorange}\@thefnmark}}
}

\makeatother

%%%%%%%%%%%%%%%%%%%%%%%%%%%%%%%%%%%%%%%%%%%%%%%%%%%%%%%%%%%%
% Now we can start to load the packages needed
%%%%%%%%%%%%%%%%%%%%%%%%%%%%%%%%%%%%%%%%%%%%%%%%%%%%%%%%%%%%

%%%%%%%%%% General      %%%%%%%%%%
% koma comes with more advanced handling of tocs, scrhack
% makes other packages like listings work with that
% TODO this might cause problems with updated packages
\usepackage{scrhack}
% we want to support both pdfLaTeX and XeLaTeX
% TODO consider luaTeX as well
\usepackage{ifxetex}
% language settings, output-encoding
% localisation (languages set in document directive):
\usepackage{babel}
% % TODO context sensitive quotes, investigate for biblatex
% \usepackage{csquotes}
% input encoding to support diacritics
\usepackage[utf8]{inputenc}
% output encoding:
\usepackage[T1]{fontenc}

%%%%%%%%%% Look & Feel  %%%%%%%%%%
% define colors as you need them
\usepackage{color}
% nicer headings/footer:
\usepackage{scrpage2}
% typeset URLs nicely:
\usepackage{url}
% extended table environment:
\usepackage{tabularx}
% % colours in tables:
% \usepackage{colortbl}
% % booktabs for nice looking tables
% \usepackage{booktabs}
% used for adapting list environment, e.g. for glossaries:
\usepackage{enumitem}
% include source code listings in the text:
\usepackage{listings}

%%%%%%%%%% Graphics     %%%%%%%%%%
% we need to be able to include graphics:
\usepackage{graphicx}
% we use colours in tables etc.:
\usepackage[dvipsnames,usenames]{xcolor}
% TikZ is cool for diagrams:
\usepackage{tikz}
% these 2 packages can be used to generate
% graphs & tables from tabular data:
\usepackage{pgfplots}
\usepackage{pgfplotstable}
% pinning to specific pgfplots version
\pgfplotsset{compat=1.8}
% make text flow around images
% \usepackage{picins}

%%%%%%%%%% Color Definitions %%%%%
% color definitions for listings (example in java)
\definecolor{pblue}{rgb}{0.13,0.13,1}
\definecolor{pgreen}{rgb}{0,0.5,0}
\definecolor{pred}{rgb}{0.9,0,0}
\definecolor{pgrey}{rgb}{0.46,0.45,0.48}

% color definitions corporate design vzg and coli-conc
\definecolor{vzgheadline}{RGB}{0,102,102}
\definecolor{vzgpetrol}{RGB}{47,140,144}
\definecolor{vzgblue}{RGB}{84,128,187}
\definecolor{coliorange}{HTML}{C14E00}
\definecolor{coliblack}{HTML}{393939}
\definecolor{coligrey}{HTML}{E3E3DF}

%%%%%%%%%% lstset Definitions %%%%%
\lstset{language=Java,
  frame=single,
  showspaces=false,
  showtabs=false,
  breaklines=true,
  showstringspaces=false,
  breakatwhitespace=true,
  captionpos=b,
  nolol=false
  commentstyle=\color{pgreen},
  keywordstyle=\color{pblue},
  stringstyle=\color{pred},
  basicstyle=\ttfamily,
  moredelim=[il][\textcolor{pgrey}]{\$\$},
  moredelim=[is][\textcolor{pgrey}]{\%\%}{\%\%}
}

%%%%%%%%%% Bib          %%%%%%%%%%
% allows us to use author-year styles:
\usepackage[round,colon,sort]{natbib}
% \usepackage{natbib}
% allows to have a separate bibliography for weblinks (or more):
\usepackage{multibib}

%%%%%%%%%% Other        %%%%%%%%%%
% % is already loaded automatically:
% \usepackage{calc}
% % do we want an index?
% \usepackage{makeidx}
% include hyperlinks in PDF and much more:
% FIXME Chose between coloured or black
\usepackage[
  colorlinks=true,
  linkcolor=black,
  menucolor=coliorange,
  citecolor=coliorange,
  urlcolor=coliorange,
]{hyperref}
\hypersetup{linktoc=all}

%%%%%%%%%%%%%%%%%%%%%%%%%%%%%%%%%%%%%%%%%%%%%%%%%%%%%%%%%%%%
% Set up header and footer
%%%%%%%%%%%%%%%%%%%%%%%%%%%%%%%%%%%%%%%%%%%%%%%%%%%%%%%%%%%%

\pagestyle{scrheadings}
\clearscrheadfoot
\renewcommand{\headfont}{\normalfont\rmfamily\itshape}
\ohead{\headmark}
% FIXME choose between english and german version
\lofoot[Coli-conc Report No. n]{Coli-conc Report No. n}
% \lofoot[Coli-conc Report Nr. n]{Coli-conc Report Nr. n}
\ofoot[\pagemark]{\pagemark}

%%%%%%%%%%%%%%%%%%%%%%%%%%%%%%%%%%%%%%%%%%%%%%%%%%%%%%%%%%%%
%   lists of... and captions
%%%%%%%%%%%%%%%%%%%%%%%%%%%%%%%%%%%%%%%%%%%%%%%%%%%%%%%%%%%%

% for listings, we define the names used in the captions
% and in the list of listings
% we also take the -verzeichnis out of the list of
% figures and list of tables
% TODO differences in defined languages/babel; Text formating
\AtBeginDocument{%
  \makeatletter
    \addto\captionsngerman{%
      \def\lstlistingname{Quelltext}%
      \def\lstlistlistingname{Quelltexte}%
%       \renewcommand{\listfigurename}{\normalfont\rmfamily Abbildungen}%
%       \renewcommand{\listtablename}{\normalfont\rmfamily Tabellen}%
      \renewcommand{\listfigurename}{Abbildungen}%
      \renewcommand{\listtablename}{Tabellen}%
    }%
  \makeatother
}
% % for other than KOMA classes, the following may have
% % to be used instead:
% \renewcommand*{\lstlistlistingname}{Quelltexte}
% \renewcommand*{\lstlistingname}{Quelltexte}

% captions should be set in a small font size
\addtokomafont{caption}{\footnotesize}
% \addtokomafont{captionlabel}{\bfseries}

%%%%%%%%%%%%%%%%%%%%%%%%%%%%%%%%%%%%%%%%%%%%%%%%%%%%%%%%%%%%
%   General changes to formatting
%%%%%%%%%%%%%%%%%%%%%%%%%%%%%%%%%%%%%%%%%%%%%%%%%%%%%%%%%%%%

% this is the KOMA-way of changing titles.
%\setkomafont{sectioning}{\rmfamily\bfseries}
\setkomafont{disposition}{\color{coliorange}\rmfamily\bfseries}
  
% we want an url to be printed in the standard font
\urlstyle{same}

%%%%%%%%%%%%%%%%%%%%%%%%%%%%%%%%%%%%%%%%%%%%%%%%%%%%%%%%%%%%
% Set up citation styles
%%%%%%%%%%%%%%%%%%%%%%%%%%%%%%%%%%%%%%%%%%%%%%%%%%%%%%%%%%%%

% we use multibib.sty to generate a weblinks bibliography,
% define more as you see fit (Software, RFCs, Images)
\newcites{web}{Weblinks}
% % FIXME If you use several bibliographies, you need to run
% % all of them through bibTeX. This can be achieved with the
% % following script (for Unix, Mac, and Linux):
% #!/bin/bash
% for file in *.aux ; do
%   bibtex `basename $file .aux`
% done

%%%%%%%%%%%%%%%%%%%%%%%%%%%%%%%%%%%%%%%%%%%%%%%%%%%%%%%%%%%%
% float control
%%%%%%%%%%%%%%%%%%%%%%%%%%%%%%%%%%%%%%%%%%%%%%%%%%%%%%%%%%%%
% more sensible defaults for how much space floats can take
% up on pages with text:
% http://www.tug.org/texmf-dist/doc/generic/FAQ-en/html/FAQ-floats.html
\setcounter{topnumber}{2}
\setcounter{bottomnumber}{9}
\setcounter{totalnumber}{20}
\setcounter{dbltopnumber}{9}
\renewcommand{\topfraction}{0.85}
\renewcommand{\bottomfraction}{0.7}
\renewcommand{\textfraction}{0.15}
\renewcommand{\floatpagefraction}{0.7}
\renewcommand{\dbltopfraction}{.7}
\renewcommand{\dblfloatpagefraction}{.7}


% FIXME in the following file all necessary input has to be made
%%%%%%%%%%%%%%%%%%%%%%%%%%%%%%%%%%%%%%%%%%%%%%%%%%%%%%%%%%%%
%   General information input part
%%%%%%%%%%%%%%%%%%%%%%%%%%%%%%%%%%%%%%%%%%%%%%%%%%%%%%%%%%%%

% FIXME enter document title
\newcommand{\colititle}{Coli-conc Report Template}

% FIXME enter the subject(s) divided by comma
\newcommand{\colisubject}{Template, VZG, Coli-conc, Report}

% FIXME real date of delivery. \today is todays date as dd. MM yyyy
\newcommand{\colidate}{\today}
% \newcommand{\colidate}{1. April 1970}

% FIXME enter the e-mail-address
\newcommand{\colimail}{coli-conc@gbv.de}


%%%%%%%%%%%%%%%%%%%%%%%%%%%%%%%%%%%%%%%%%%%%%%%%%%%%%%%%%%%%
%   For coli-documentation and coli-report only
%%%%%%%%%%%%%%%%%%%%%%%%%%%%%%%%%%%%%%%%%%%%%%%%%%%%%%%%%%%%

% FIXME (un)mark the version that is needed
% \newcommand{\colireportlang}{No. }
\newcommand{\colireportlang}{Nr. }

% FIXME enter the report number
\newcommand{\colireportno}{11}


%%%%%%%%%%%%%%%%%%%%%%%%%%%%%%%%%%%%%%%%%%%%%%%%%%%%%%%%%%%%
%   For coli-article and coli-report only
%%%%%%%%%%%%%%%%%%%%%%%%%%%%%%%%%%%%%%%%%%%%%%%%%%%%%%%%%%%%

% Number of authors and working addresses has to correspond.
% The tex-file has to be adjusted accordingly.
% FIXME enter the document author(s)
\newcommand{\coliauthorone}{Mustermann, Max}
\newcommand{\coliauthortwo}{Musterfrau, Martha}
\newcommand{\coliauthorthree}{Bla}
\newcommand{\coliauthorfour}{Blubb}


%%%%%%%%%%%%%%%%%%%%%%%%%%%%%%%%%%%%%%%%%%%%%%%%%%%%%%%%%%%%
%   For coli-article only
%%%%%%%%%%%%%%%%%%%%%%%%%%%%%%%%%%%%%%%%%%%%%%%%%%%%%%%%%%%%

% FIXME enter the author(s) working address(es)
% has to be the same amount as authors in the part before!
\newcommand{\coliaddressone}{Verbundzentrale des GBV (VZG), Göttingen}
\newcommand{\coliaddresstwo}{Department of Computer Science, University of Wherever}
\newcommand{\coliaddressthree}{Department of Biology, University of Wherever}
\newcommand{\coliaddressfour}{Department of Dingsda, University of Wherever}


%%%%%%%%%%%%%%%%%%%%%%%%%%%%%%%%%%%%%%%%%%%%%%%%%%%%%%%%%%%%
%   For coli-documentation only
%%%%%%%%%%%%%%%%%%%%%%%%%%%%%%%%%%%%%%%%%%%%%%%%%%%%%%%%%%%%

% FIXME enter the document author(s)
\newcommand{\coliauthor}{Mustermann, Max und Musterfrau, Martha}

% FIXME (un)mark the needed text type
% \newcommand{\colitexttype}{Technical Documentation}
\newcommand{\colitexttype}{Technische Dokumentation}

% FIXME (un)mark what you need/want
% text type + enumeration
% \newcommand{\colidokfooter}{\colitexttype~\colireportlang~\colireportno}
% only text type
% \newcommand{\colidokfooter}{\colitexttype}
% only title
\newcommand{\colidokfooter}{\colititle}


%%%%%%%%%%%%%%%%%%%%%%%%%%%%%%%%%%%%%%%%%%%%%%%%%%%%%%%%%%%%
%   For coli-report only
%%%%%%%%%%%%%%%%%%%%%%%%%%%%%%%%%%%%%%%%%%%%%%%%%%%%%%%%%%%%

% FIXME enter the author(s) e-mail-address(es)
% has to be the same amount as authors specified in the section
% For coli-article and coli-report only
\newcommand{\colimailone}{max.mustermann@vzg.de}
\newcommand{\colimailtwo}{martha.musterfrau@vz.de}
\newcommand{\colimailthree}{bla@vzg.de}
\newcommand{\colimailfour}{blubb@vzg.de}

% FIXME enter the affiliation
\newcommand{\coliinstitute}{Verbundzentrale des GBV (VZG), Göttingen}

% FIXME enter the DOI address
\newcommand{\colidoi}{<DOI Address>}

% FIXME enter the used licence
\newcommand{\colilicence}{CC-BY-SA}

% nothing to enter here!
\newcommand{\colifooter}{Coli-conc Report\colireportlang~\colireportno}


%%%%%%%%%%%%%%%%%%%%%%%%%%%%%%%%%%%%%%%%%%%%%%%%%%%%%%%%%%%%
% lstset Definitions
%%%%%%%%%%%%%%%%%%%%%%%%%%%%%%%%%%%%%%%%%%%%%%%%%%%%%%%%%%%%
% FIXME define for the language you are really using
% Java is just an example, but usable if needed
\lstset{language=Java,
  frame=single,
  showspaces=false,
  showtabs=false,
  breaklines=true,
  showstringspaces=false,
  breakatwhitespace=true,
  captionpos=b,
  nolol=false,
  numbers=left,
  numbersep=5pt,
  numberstyle=\tiny\color{pgrey},
  commentstyle=\color{pgreen},
  keywordstyle=\color{pblue},
  stringstyle=\color{pred},
  basicstyle=\ttfamily,
  moredelim=[il][\textcolor{pgrey}]{\$\$},
  moredelim=[is][\textcolor{pgrey}]{\%\%}{\%\%}
}


% do we want to include a glossary and an index?

% \makeglossary
\makeindex

%%%%%%%%%%%%%%%%%%%%%%%%%%%%%%%%%%%%%%%%%%%%%%%%%%%%%%%%%%%%
% Begin Real Document
%%%%%%%%%%%%%%%%%%%%%%%%%%%%%%%%%%%%%%%%%%%%%%%%%%%%%%%%%%%%
\begin{document}

% stuff before real content, no chapter marks, roman page numbers
\frontmatter
\pagenumbering{Roman}

% % the cite commands will generate entries in the index, if used
% \citeindextrue

%%%%%%%%%%%%%%%%%%%%%%%%%%%%%%%%%%%%%%%%%%%%%%%%%%%%%%%%%%%%
% We don't use a custom \maketitle for now
% so we have to make our own title page
%%%%%%%%%%%%%%%%%%%%%%%%%%%%%%%%%%%%%%%%%%%%%%%%%%%%%%%%%%%%

\begin{titlepage}
  \hypersetup{%
    pdftitle={\colititle},
    pdfauthor={\coliauthor},
    pdfsubject={\colisubject}
  }
 
\par\medskip
\noindent\makebox[\linewidth][c]{%
	\colorbox{vzgblue}{%
    	\parbox[l][1.2em][c]{\paperwidth}{%
      		\hspace*{\dimexpr\hoffset+\oddsidemargin+1in\relax}%
		    \begin{minipage}{.3\textwidth}
        		\includegraphics[width=\linewidth, height=1.7em, keepaspectratio]{figures/vzg.png}
		    \end{minipage}%
		    \begin{minipage}{.05\textwidth}
        		\mbox{}
		    \end{minipage}%
		    \begin{minipage}{.7\textwidth}
			    \begin{flushright}
			    	\textbf{{\footnotesize \textcolor{white}{Verbundzentrale des GBV}}}\relax
			    \end{flushright}
			\end{minipage}%
    	}%
	}%
}\par\medskip
  
  \begin{center}
%	\colorbox{coligrey}{\includegraphics[scale=0.15]{figures/vzg.png}}
%    \includegraphics[scale=0.5]{figures/vzg-schriftzug.png}\\

    \vfill

    {
%	FIXME select used language    
      \selectlanguage{english} % switching to english
      % \selectlanguage{ngerman} % switching to german

      \Huge
      \bfseries
      \textcolor{vzgblue}{
	    \colitexttype
		}
		
	\Large
	\bfseries
	\textcolor{vzgblue}{
	Report \colireportlang \colireportno
	}
      
      \vspace{1.25cm}
      
      \huge
      \mdseries
        \colititle
    }

    \vfill

    {
      \normalsize

      \vfill
      \vfill
      
      Göttingen, \colidate
    }
      
      \vfill
      
    {
    \includegraphics[scale=0.3]{figures/coli-conc-wave.png}
    %{\color{vzgblue}\rule[0cm]{\textwidth}{2cm}}\\
    %\colorbox{vzgblue}{\includegraphics[scale=0.7]{figures/coli-conc-wave.png}}
%	  \includegraphics[scale=0.7]{figures/vzg-schriftzug.png}
%		\footnotesize 
%		\selectlanguage{ngerman}
%        Verbundzentrale des GBV (VZG) $\cdot$ Platz der Göttinger Sieben 1 $\cdot$ 37073 Göttingen
    }
  \end{center}
\end{titlepage}


% if using tocloft.sty, we need an explicit clearpage
\clearpage
% we want chapters and sections in the toc, then generate it
\setcounter{tocdepth}{1}

\tableofcontents

% now we have real content, with section marks and egyptian numerals
\mainmatter

%%%%%%%%%%%%%%%%%%%%%%%%%%%%%%%%%%%%%%%%%%%%%%%%%%%%%%%%%%%%
% WHERE THE REAL WORK TAKES PLACE
%%%%%%%%%%%%%%%%%%%%%%%%%%%%%%%%%%%%%%%%%%%%%%%%%%%%%%%%%%%%
\selectlanguage{english} % switching to english
% \selectlanguage{ngerman} % switching to german

\include{text}


%%%%%%%%%%%%%%%%%%%%%%%%%%%%%%%%%%%%%%%%%%%%%%%%%%%%%%%%%%%%
% list of figures, tables
%%%%%%%%%%%%%%%%%%%%%%%%%%%%%%%%%%%%%%%%%%%%%%%%%%%%%%%%%%%%

% here we have all the stuff where chapters have no numbers etc.
% you will find a lot of \cleardoublepage and \phantomsection
% commands, these help hyperref.sty to find the right targets
% for hyperlinks
\backmatter

% The list of figures
  \cleardoublepage
  \phantomsection
%   \addcontentsline{toc}{chapter}{Abbildungen}
  \listoffigures

% The list of tables
  \cleardoublepage
  \phantomsection
%   \addcontentsline{toc}{chapter}{Tabellen}
  \listoftables

% list of listings
  \cleardoublepage
  \phantomsection
%   \addcontentsline{toc}{chapter}{Quelltexte}
  \lstlistoflistings  % generated by listings.sty

%%%%%%%%%%%%%%%%%%%%%%%%%%%%%%%%%%%%%%%%%%%%%%%%%%%%%%%%%%%%
%%%%%%%%%%%%%%%%%%%%%%%%%%%%%%%%%%%%%%%%%%%%%%%%%%%%%%%%%%%%
\chapter{Quellen}
% \addcontentsline{toc}{chapter}{Quellen}
%%%%%%%%%%%%%%%%%%%%%%%%%%%%%%%%%%%%%%%%%%%%%%%%%%%%%%%%%%%%

% While it is possible to cite something without putting
% in a reference, the question is, why would you do that?
% \nociteweb{imai}

\phantomsection
\renewcommand{\bibname}{Literatur}
\bibliography{bibliography}
\bibliographystyle{abbrvnat}
\addcontentsline{toc}{section}{Literatur}

\phantomsection
\bibliographystyleweb{imai}
\bibliographyweb{bibliography}
\addcontentsline{toc}{section}{Weblinks}


\section*{Software}
\addcontentsline{toc}{section}{Software}

\LaTeX, 08. März 2012, \url{http://www.latex-project.org/}

openSUSE 12.1, 08. März 2012, \url{http://www.opensuse.org/}

%%%%%%%%%%%%%%%%%%%%%%%%%%%%%%%%%%%%%%%%%%%%%%%%%%%%%%%%%%%%
% we now go into the appendix. basically the appendix is one chapter
% with some sections below. The appendix chapter has no mark, the
% sections alphabetical ones, subsections will have numerical ones
% TODO format lower level sections in appendix
\renewcommand{\thechapter}{\Alph{chapter}}
\renewcommand{\thesection}{\Alph{section}}
\renewcommand{\thesubsection}{\Alph{section}.\arabic{subsection}}
\setcounter{section}{0} % need to be explicit, since they are not reset in back matter
\addtocounter{chapter}{1} % help hyperref find the correct section

%%%%%%%%%%%%%%%%%%%%%%%%%%%%%%%%%%%%%%%%%%%%%%%%%%%%%%%%%%%%
%%%%%%%%%%%%%%%%%%%%%%%%%%%%%%%%%%%%%%%%%%%%%%%%%%%%%%%%%%%%
\chapter{Anhänge}\label{chapter:appendix}
% \addcontentsline{toc}{chapter}{Anhänge}
%%%%%%%%%%%%%%%%%%%%%%%%%%%%%%%%%%%%%%%%%%%%%%%%%%%%%%%%%%%%

%  Umfangreiche zusätzliche Informationen, die im Textverlauf stören würden, die aber für 
%  die Arbeit wichtig sind, wie z.B. Programmcode, Fragebögen, Evaluationstabellen

%%%%%%%%%%%%%%%%%%%%%%%%%%%%%%%%%%%%%%%%%%%%%%%%%%%%%%%%%%%%
\section{Programmcode}\label{sec:appendix_code}
%%%%%%%%%%%%%%%%%%%%%%%%%%%%%%%%%%%%%%%%%%%%%%%%%%%%%%%%%%%%

%  Es geht hier nicht darum, den kompletten Quelltext des praktischen Teils der Arbeit 
%  abzudrucken
  
%  Es sollen lediglich besonders wichtige Fragmente (API-Definitionen, 
%  Kommunikationsprotokolle) dokumentiert werden


\end{document}